\documentclass{report}
\usepackage{amsmath,amssymb}
\setlength{\parindent}{0mm}
\setlength{\parskip}{1em}
\begin{document}
\begin{center}
\rule{6in}{1pt} \
{\large Philip Browne \\
{\bf Convergence of the ESO method for topology optimization}}

Department of Mathematical Sciences \\ University of Bath \\ Bath \\ UK \\ BA2 7AY
\\
{\tt pab23@bath.ac.uk}\\
Chris Budd\end{center}

Evolutionary structural optimization (ESO) is a method for topology
optimization developed by engineers in the early 1990s. Whilst there
exists a great deal of \textit{engineering intuition} about the
algorithm, its mathematical foundation has repeatedly come into question.

This work seeks to explain the non-monotonic behaviour of the ESO
algorithm by showing how the nonlinearities inherent in the underlying
state equations leads to behaviour not predicted by the infinitesimal
information used by the ESO algorithm.

Furthermore, we show how local h-refinement of the finite element mesh
can show descent of the algorithm.


\end{document}
